% Abschnitt 5.1 Handlungsbedarf – ENTWURF AP4 (17.11.2026, Claude). Plan: review/plan-kap5.md, Abschnitt 5 „5.1 Handlungsbedarf“.
% Belege nur mit Bewertung „trägt“ oder „trägt mit Einschränkung“ aus review/kap5-entwurf/belegpruefung.md (AP2).
% Wird in AP7 unter \section{Handlungsbedarf}\label{sec:handlungsbedarf} in Kapitel/05_Handlungsbedarf_Methodik.tex eingefügt.
% AP7: Verweis auf Tabelle~\ref{tab:stacks} (Merkmalsvergleich in Abschnitt 4.4) ergänzen, sobald die Tabelle dort steht
% (Plan Abschnitt 9, Abhängigkeit D2); bis dahin tragen die Belege im Fließtext die Wertung.
% Kann-Tabelle tab:defizite (AP3): falls übernommen, im letzten Absatz vor der Tabelle mit \ref einführen.

% Basis: Zielsetzung und Forschungsfragen F1–F4 laut Gerüst Kapitel/01_Einleitung.tex (Abschnitt 1.2);
% Betrachtungsumfang: in Kap. 3 und 4 dargestellte Arbeiten, Recherche Stand November 2026 (review/plan-sdt-verwandte-arbeiten.md).
Maßstab der folgenden Wertung ist das Ziel der Arbeit aus Abschnitt~\ref{sec:zielsetzung}: eine aufgabenunabhängige Umgebung, in der Anwender ohne Robotikerfahrung den Nutzen eines Humanoiden prüfen, indem sie Funktionen des Unitree G1 zuerst in der Simulation und anschließend am realen Roboter erproben.
Ein solcher Zugang setzt an Potenzial und Hemmnissen humanoider Roboter an.
Laut einer Übersichtsarbeit können Humanoide in der Industrie repetitive und arbeitsintensive Aufgaben übernehmen und mit Menschen bei Montage, Verpackung und Materialhandhabung zusammenarbeiten \cite{Tong2024}\ev{Tong2024.08}, während hohe Entwicklungs- und Wartungskosten ihre kommerzielle Verbreitung begrenzen \cite{Tong2024}\ev{Tong2024.09}.
Kommerziell verfügbare Humanoid-Hardware ist zudem nach Einschätzung einer Arbeit zu einem quelloffenen Humanoiden überwiegend teuer, nicht quelloffen und intransparent \cite{Chi2025}\ev{Chi2025.01}.
Für Humanoide wie Pepper, die unter Verweis auf weitere Arbeiten als schwer und teuer zu bauen und zu testen gelten, werden Simulationsumgebungen als entscheidend angesehen, um neue Konzepte sicher und wirksam zu erproben \cite{Sekkat2024}\ev{Sekkat2024.05}.
\qf{Hemmnisse beim Einstieg in humanoide Robotik aus Sicht von Unternehmen (Kosten, Risiko, fehlendes Fachwissen) – Acatech2020 und FAPS2024 sind freigegeben, haben aber kein PDF}
Aus dem Ziel folgen vier Gesichtspunkte der Wertung, einer je Forschungsfrage: die Durchgängigkeit von der Simulation zum Roboter (F1), die Aussagekraft der Simulation (F2), Bedienung und Sicherheit für Personen ohne Robotikerfahrung (F3) sowie die Offenheit für vorab nicht bekannte Aufgaben (F4).
Gewertet werden die in den Kapiteln~\ref{chap:simulation} und~\ref{chap:faehigkeiten} dargestellten Arbeiten und Werkzeuge nach dem Stand der Recherche im November~2026; die daraus abgeleiteten Defizite D1 bis D4 gelten nur für diesen Betrachtungsumfang.
%
% Basis: Kap. 4.4 (sec:verwandte_arbeiten), Belege L2–L5 in review/kap5-entwurf/belegpruefung.md;
% Negativbefund G1Pilot: Suche im README-Text vom 09.11.2026 (Commit 5346e52) ohne Treffer für sim/gazebo/mujoco (AP2, HucebotG1Pilot.05);
% Negativbefund OpenWBT: README (Commit a0154ca) nennt Simulation, aber nicht, ob derselbe Programmstand läuft.
Das erste Defizit~(D1) betrifft die Durchgängigkeit: Unter den betrachteten Arbeiten deckt eine gemeinsame Ausführung in Simulation und Realbetrieb jeweils nur einen Teil des Funktionsumfangs des G1 ab.
Der Simulator unitree\_mujoco des Herstellers bindet nach dessen Angabe mit dem Hersteller-SDK entwickelte Steuerprogramme ein und soll so einen nahtlosen Übergang von der Simulation zur Entwicklung am Roboter ermöglichen \cite{UnitreeMujoco}\ev{UnitreeMujoco.01}.
In der ausgewerteten Fassung der Projektbeschreibung unterstützt er jedoch nur die Entwicklung auf niedriger Ebene und dient vor allem der Sim-to-Real-Prüfung von Reglern \cite{UnitreeMujoco}\ev{UnitreeMujoco.02}.
Der Onboard-Regler, der am realen G1 die Fortbewegung regelt (Abschnitt~\ref{sec:fortbewegung_sdt}), ist dagegen nicht quelloffen, sodass eine Vergleichsarbeit ihm keine inneren Mechanismen zuschreibt \cite{Hu2026}\ev{Hu2026.02}.
Ein Ganzkörper-Teleoperationssystem für den G1, beschrieben in einem nicht begutachteten Preprint, überträgt die mit einem IMU-basierten Motion-Capture-Anzug aufgenommene Bewegung auf den Roboter \cite{Durrani2026}\ev{Durrani2026.01}.
Es wird zuerst mit dem MuJoCo-Modell des G1 geprüft und danach laut den Autoren ohne Änderung am realen Roboter eingesetzt \cite{Durrani2026}\ev{Durrani2026.02}, wobei derselbe Retargeting-Code unverändert in beiden Umgebungen läuft \cite{Durrani2026}\ev{Durrani2026.03} und die Stabilität auf niedriger Ebene den Onboard-Reglern des G1 überlassen bleibt \cite{Durrani2026}\ev{Durrani2026.04}.
OpenWBT unterstützt laut Projektbeschreibung reale Roboter ebenso wie Simulationsumgebungen und steuert dabei den G1 über eine Apple Vision Pro \cite{GalbotOpenWBT}\ev{GalbotOpenWBT.01}; ob in beiden Fällen derselbe Programmstand ausgeführt wird, geht aus der Beschreibung nicht hervor.
Andere Arbeiten bleiben in einer der beiden Umgebungen.
Eine Studie zum G1 simuliert ausschließlich in MATLAB/Simulink mit einem selbst erstellten Robotermodell \cite{Yu2026}\ev{Yu2026.03} und nennt Versuche am realen G1 zur Verifikation der Simulationsergebnisse als notwendig \cite{Yu2026}\ev{Yu2026.04} und als künftige Arbeit \cite{Yu2026}\ev{Yu2026.11}.
G1Pilot überlässt den Unterkörper dem Regler des Herstellers und stellt Werkzeuge für die Steuerung des Oberkörpers und die Teleoperation bereit \cite{HucebotG1Pilot}\ev{HucebotG1Pilot.01}; die Projektbeschreibung nennt eine Visualisierung des realen Roboters \cite{HucebotG1Pilot}\ev{HucebotG1Pilot.05}, einen Simulationsbetrieb beschreibt sie nicht.
Auch eine Übersichtsarbeit hält fest, dass trotz der Fortschritte in der Simulation der sichere Übergang von der Simulation in die reale Welt eine Herausforderung bleibt \cite{Tong2024}\ev{Tong2024.06}.
Aus D1 folgt die Forschungsfrage F1 (Abschnitt~\ref{sec:zielsetzung}).
%
% Basis: Kap. 3.1.2 (sec:sim2real), 3.1.3 (sec:abbildungsguete), Belege L6–L7 in review/kap5-entwurf/belegpruefung.md;
% Wertung „keine Bewertung für einen Bedien-Stack“: keine der in Kap. 3.1.3 und 4.4 betrachteten Arbeiten vergleicht Simulation und
% Realbetrieb für eine Umgebung mit unterschiedlichen Reglern für die Fortbewegung (eigene Auswertung, review/plan-sdt-verwandte-arbeiten.md).
Das zweite Defizit~(D2) betrifft die Aussagekraft der Simulation: Die betrachtete Literatur bewertet sie für gelernte Regelungsstrategien und einzelne Systeme, nicht für eine Bedienumgebung als Ganzes.
Abschnitt~\ref{sec:sim2real} beschreibt die Ursachen, aus denen Simulation und realer Roboter voneinander abweichen.
Im Zusammenhang mit dem Transfer gelernter Ganzkörperbewegungen auf den G1 ist die Sim-to-Real-Lücke laut den Autoren ausgeprägter als die Abweichung zwischen zwei Simulatoren, unter anderem wegen verrauschter Sensoreingaben, ungenauer Robotermodelle und Unterschieden der Aktoren \cite{He2025}\ev{He2025.07}.
Für Manipulatoren gilt die Simulation von kinematischem Modell und Regelung nach einem Vergleich aus dem Jahr~2019 als weitgehend gelöst, für Interaktionen zwischen simulierten Objekten sind dagegen erhebliche Weiterentwicklungen nötig \cite{Collins2019}\ev{Collins2019.09}.
Nach einer Arbeit zu Pepper, die sich dabei auf weitere Literatur stützt, sind Simulationsumgebungen kostengünstig und sicher, bilden die komplexe Dynamik realer Interaktionen aber nur schwer genau ab \cite{Sekkat2024}\ev{Sekkat2024.09}.
Physisches Testen vor dem Einsatz bleibt laut einer Übersichtsarbeit zwingend \cite{Choi2021}\ev{Choi2021.02}, und dieselbe Arbeit sieht Bedarf an neuen Wegen zur Validierung von Modellen, besonders für noch nicht gebaute Roboter \cite{Choi2021}\ev{Choi2021.18}.
Die in Abschnitt~\ref{sec:abbildungsguete} vorgestellten Bewertungswege unterscheiden sich im Bezug: Eine Arbeit zur Navigation misst einen Simulator primär nicht an seinem visuellen oder physikalischen Realismus, sondern an seiner Vorhersagekraft, also daran, ob die Rangfolge zweier Methoden aus der Simulation in der Realität erhalten bleibt \cite{Kadian2020}\ev{Kadian2020.01}.
Das Teleoperationssystem für den G1 prüft in der Simulation die Korrektheit des Retargetings, problematische Gelenkkonfigurationen und die Stabilität der getesteten Bewegungen \cite{Durrani2026}\ev{Durrani2026.05} und wertet eine bestandene Prüfung, weil derselbe Retargeting-Code in beiden Umgebungen läuft, als starkes Indiz für die Übertragbarkeit \cite{Durrani2026}\ev{Durrani2026.03}.
Für eine Umgebung, die am realen G1 den nicht quelloffenen Onboard-Regler nutzt und in der Simulation folglich einen anderen Regler für die Fortbewegung einsetzen muss, liegt unter den betrachteten Arbeiten keine Bewertung der Abbildungsgüte vor.
Aus D2 folgt die Forschungsfrage F2 (Abschnitt~\ref{sec:zielsetzung}).
%
% Basis: Kap. 4.3.1 (sec:teleoperation_teil), 4.3.2 (sec:laienbedienung_sdt), 3.4.1 (sec:normen), 3.4.3 (sec:sicherheit_zweibeiner);
% Belege L8–L11 in review/kap5-entwurf/belegpruefung.md. Wertung „Werkzeug für Personen ohne Programmiererfahrung nur für Nao“:
% unter den in Kap. 4.3.2 betrachteten Arbeiten ist Choregraphe das einzige humanoid-spezifische Werkzeug (eigene Auswertung).
Das dritte Defizit~(D3) betrifft Bedienung und Sicherheit: Die betrachteten Bedienkonzepte für den G1 setzen Spezialhardware oder Vorwissen voraus, ein Werkzeug für Personen ohne Programmiererfahrung liegt nur für einen anderen Humanoiden vor, und der bestehende Normenrahmen ist nicht auf zweibeinige Roboter ausgelegt.
Frühere Arbeiten zur Teleoperation von Humanoiden nutzen nach der Darstellung einer kamerabasierten Arbeit Motion-Capture-Systeme, VR-Brillen mit Ich-Perspektive und Exoskelette, die teuer und an bestimmte Einsatzorte gebunden sind \cite{Fu2024}\ev{Fu2024.01}; für die Datensammlung in großem Maßstab sind solche Aufbauten laut einer weiteren Arbeit hinderlich \cite{He2024}\ev{He2024.03}, und bisherige Ansätze der Teleoperation zur Datensammlung leiden unter begrenzter Skalierbarkeit und aufwendiger Einrichtung \cite{Zhao2025}\ev{Zhao2025.01} (Abschnitt~\ref{sec:teleoperation_teil}).
Für den G1 steuert die Teleoperation des Herstellers den Roboter über XR-Brillen wie Apple Vision Pro, PICO~4 Ultra Enterprise oder Meta Quest~3 \cite{UnitreeXRTeleop}\ev{UnitreeXRTeleop.01}, und Neueinsteigern empfiehlt die Projektbeschreibung, vorher die offizielle Dokumentation mindestens bis zum Kapitel zur Anwendungsentwicklung zu lesen \cite{UnitreeXRTeleop}\ev{UnitreeXRTeleop.02}.
Ein weiteres, als Preprint veröffentlichtes System kombiniert für Bauaufgaben mit dem G1 eine XR-Steuerung des Oberkörpers mit Pedalen für die Fortbewegung \cite{AliPour2026}\ev{AliPour2026.01}.
Werkzeuge für Nutzer ohne Erfahrung in der Roboterprogrammierung verfolgt die End-User-Roboterprogrammierung, deren Ziel es ist, die Komplexität der Roboterprogrammierung in auch ohne umfangreiche Robotikerfahrung zugängliche Methoden zu überführen \cite{Ajaykumar2021}\ev{Ajaykumar2021.03} (Abschnitt~\ref{sec:laienbedienung_sdt}).
Für einen Humanoiden ist Choregraphe ein solches Werkzeug, eine grafische Programmierumgebung für den Roboter Nao \cite{Pot2009}\ev{Pot2009.01}, die nach Absicht der Entwickler Generalisten die schnelle Entwicklung komplexer Anwendungen ermöglichen soll \cite{Pot2009}\ev{Pot2009.02} und laut den Autoren die Arbeit mit dem Roboter erlaubt, ohne eine Zeile C++ zu kompilieren \cite{Pot2009}\ev{Pot2009.06}.
Unter den betrachteten Arbeiten richtet sich ein solches Werkzeug damit an den Nao, nicht an den G1; zudem zeigen akademische Arbeiten zur End-User-Roboterprogrammierung nach einem Survey aus dem Jahr~2021 wegen fehlender Einsätze und Nutzerstudien außerhalb des Labors möglicherweise nur ein eingeschränktes Bild \cite{Ajaykumar2021}\ev{Ajaykumar2021.12}.
Für den sicheren Betrieb setzen die meisten Robotiknormen laut einem nicht begutachteten Bericht der IEEE Humanoid Study Group feste oder statisch stabile Systeme voraus \cite{Prather2025}\ev{Prather2025.01} und lassen sich nicht direkt auf Humanoide anwenden \cite{Prather2025}\ev{Prather2025.06} (Abschnitt~\ref{sec:normen}).
Selbst mit gebremsten und abgeschalteten Gelenken befindet sich ein Humanoid laut diesem Bericht nicht notwendig in einem sicheren Zustand \cite{Prather2025}\ev{Prather2025.03}, und nach einem nicht begutachteten Preprint erzeugt der Energieentzug bei einem laufenden Zweibeiner einen unkontrollierten Sturz, der selbst eine Gefährdung ist \cite{Ding2026}\ev{Ding2026.03} (Abschnitt~\ref{sec:sicherheit_zweibeiner}).
Die Bordrechner des G1 sind laut demselben Preprint handelsübliche Standardhardware, sodass der in der dortigen Fallstudie darauf betriebene Sicherheitscontroller keine zertifizierte Sicherheitslaufzeit ist \cite{Ding2026}\ev{Ding2026.05}.
Für die Policy des Teleoperationssystems OmniH2O nennen die Autoren als Grenze fehlende Garantien und Sicherheitsprüfungen für extreme Störungen und Bewegungsziele außerhalb der Trainingsverteilung \cite{He2024}\ev{He2024.04}, und eine Übersichtsarbeit nennt die sichere Nutzung humanoider Roboter eine große Herausforderung \cite{Tong2024}\ev{Tong2024.14}.
Aus D3 folgt die Forschungsfrage F3 (Abschnitt~\ref{sec:zielsetzung}).
%
% Basis: Kap. 4.1 (tab:policies), 4.3.1, 4.4, 3.3.1 (sec:ros2); Belege L12–L13 in review/kap5-entwurf/belegpruefung.md.
% Wertung „keine gemeinsame Grundlage zum Andocken eigener Aufgaben“: kein Beleg möglich (L12 teilweise), eigene Auswertung der in
% Kap. 4 betrachteten Werkzeuge (review/plan-sdt-verwandte-arbeiten.md).
Das vierte Defizit~(D4) betrifft die Zweckbindung: Die betrachteten Werkzeuge sind jeweils auf einen Zweck zugeschnitten: Datensammlung, Training von Reglern oder Teleoperation.
Teleoperation dient in der Forschung als Methode, um Daten am Roboter für das Lernen aus Demonstrationen zu sammeln \cite{Cheng2024}\ev{Cheng2024.01}, und gilt als vorherrschende Methode der Datensammlung \cite{Zhao2025}\ev{Zhao2025.01}.
Das Projekt unitree\_rl\_mjlab ist ein Projekt für Reinforcement Learning mit MuJoCo als Physik-Backend, das unter anderem den G1 unterstützt \cite{UnitreeRLMjlab}\ev{UnitreeRLMjlab.01}, mit dem Ablauf Training, Abspielen und Übertragung auf den Roboter \cite{UnitreeRLMjlab}\ev{UnitreeRLMjlab.06}.
GR00T-WholeBodyControl stellt Checkpoints und Skripte für Training, Bewertung und Einsatz von Ganzkörperreglern für Humanoide bereit \cite{NvidiaGR00TWBC}\ev{NvidiaGR00TWBC.01} (Tabelle~\ref{tab:policies}).
Die meisten Teleoperationssysteme für Humanoide beschränken sich zudem auf Fortbewegung oder Manipulation statt auf koordinierte Ganzkörperbewegungen \cite{Ze2025}\ev{Ze2025.01}.
Einen anderen Ansatz verfolgt die Middleware ROS~2 (Abschnitt~\ref{sec:ros2}): Sie setzt Modularität als Designprinzip auf mehreren Ebenen um, von Bibliotheksschnittstellen über Nachrichtendefinitionen bis zum Software-Ökosystem \cite{Macenski2022}\ev{Macenski2022.05}, und macht seine Middleware-Implementierungen ohne Codeänderung austauschbar \cite{Macenski2022}\ev{Macenski2022.11}.
ROS wurde laut einer Übersichtsarbeit zu Humanoiden eingeführt, um die Vielfalt an Programmiersprachen und Plattformen zu vereinheitlichen und die Umsetzung zu beschleunigen \cite{Saeedvand2019}\ev{Saeedvand2019.13}.
Unter den betrachteten Werkzeugen für den G1 ist eine gemeinsame Grundlage, an die ein Anwender eigene, vorab nicht bekannte Aufgaben und Folgeprojekte andocken kann, nicht Gegenstand; die Werkzeuge verfolgen andere Ziele.
Aus D4 folgt die Forschungsfrage F4 (Abschnitt~\ref{sec:zielsetzung}).
%
% Basis: Zusammenführung D1–D4 (dieser Abschnitt); Eingrenzung laut Gerüst Kapitel/01_Einleitung.tex (Abschnitt 1.2);
% Nutzerstudie: Frage N3 in review/plan-kap5.md, Stand 16.11.2026 offen.
Zusammengefasst verbindet keine der in den Kapiteln~\ref{chap:simulation} und~\ref{chap:faehigkeiten} betrachteten Arbeiten den durchgängigen Weg von der Simulation zum realen G1 mit einer Bedienung für Personen ohne Robotikerfahrung, Schutzmechanismen für den Realbetrieb eines Zweibeiners und einer offenen Grundlage für eigene Aufgaben.
Die betrachteten Arbeiten verfolgen jeweils andere Ziele; die Lücke entsteht erst aus dem Anspruch, die vier Gesichtspunkte in einer Umgebung zu vereinen.
Daraus präzisiert sich das Ziel der Arbeit: eine Umgebung für den G1, die denselben Programmstand in Simulation und Realbetrieb ausführt, deren Abbildungsgüte für die bereitgestellten Funktionen bewertet ist, die ohne Robotikerfahrung bedienbar ist, im Realbetrieb Schutzmechanismen bereitstellt und für neue Aufgaben erweiterbar ist.
Relevant ist dieses Ziel für Anwender, die den Nutzen eines Humanoiden vor dem Einsatz erproben wollen, und für Folgeprojekte am Lehrstuhl, die auf der Umgebung aufbauen.
Nicht geschlossen werden die Defizite, die die Regelung selbst betreffen: Die Arbeit trainiert keine eigene Regelungsstrategie und überträgt keine auf den Roboter, weil D1 bis D4 die Umgebung um vorhandene Regler betreffen und nicht deren Entwurf; die Navigation am realen Roboter ist ebenfalls nicht Gegenstand (Abschnitt~\ref{sec:zielsetzung}).
Ebenso ist eine sicherheitstechnische Bewertung nicht Gegenstand: Die Arbeit stellt Schutzmechanismen für den Versuchsbetrieb bereit, keinen Sicherheitsnachweis, und sie misst die Gebrauchstauglichkeit nicht mit einer Nutzerstudie, sondern über den Bedienaufwand. \pruefen{Nutzerstudie ja/nein (N3) – bei „ja“ diesen Satz anpassen}
Welche Eigenschaften die Umgebung dafür im Einzelnen aufweisen muss, legt Abschnitt~\ref{sec:anforderungen} als prüfbare Anforderungen fest.
